% !PRINTED 468
\weberpage{468}

Nach den Formeln \S.\,115, (14) ergiebt sich zunächst sehr einfach:

\[
\mu = \frac{2\varepsilon}{7},
\]

und für $\lambda$ erhält man, wenn man die Werthe \S.\,115, (11) für $\alpha, \beta, \gamma$ einsetzt, und die Formeln \S.\,115, (12), (13) benutzt:

\weq{5}{\begin{aligned}
\lambda &= -\varepsilon\,\frac{1+\sqrt{-7}}{14} \\
\frac{\mu}{\lambda} &= -\frac{1-\sqrt{-7}}{2}.
\end{aligned}}

Setzt man daher

\[
A_1^2+A_2^2+A_3^2 = \lambda z,
\]

so ergiebt sich

\weq{6}{z = x_1^2+x_2^2+x_3^2 - \frac{1-\sqrt{-7}}{2}(x_2x_3+x_3x_1+x_1x_2),}

und diese Function wollen wir als Wurzel der Resolvente 7\textsuperscript{ten} Grades einführen. Die übrigen Wurzeln erhält man daraus durch die Substitutionen $\tau^r$, so dass sie alle in der gemeinschaftlichen Form

\weq{7}{\begin{aligned}
z_r &= \varepsilon^{2r}x_1^2 + \varepsilon^{4r}x_2^2 + \varepsilon^{r}x_3^2 \\
&\quad - \frac{1-\sqrt{-7}}{2}(\varepsilon^{-r}x_2x_3+\varepsilon^{-2r}x_3x_1+\varepsilon^{-4r}x_1x_2) \\
&\qquad r=0, 1, 2, 3, 4, 5, 6
\end{aligned}}

enthalten sind.

Die Coëfficienten der Gleichung 7\textsuperscript{ten} Grades, deren Wurzeln die sieben Grössen $z_r$ sind, sind Invarianten unserer Gruppe, deren Grad sich leicht angeben l\"asst.

Wenn nämlich $a_\nu$ der Coëfficient von $z^{7-\nu}$ in dieser Gleichung ist, nachdem der Coëfficient der siebenten Potenz auf 1 gebracht ist, so ist $a_\nu$ eine Invariante $2\nu^{\text{ten}}$ Grades. Es muss also zunächst $a_1=0$ sein, weil es keine quadratische Invariante giebt. Die übrigen Coëfficienten sind durch folgende Invariantenverbindungen linear und homogen ausgedrückt:

\weq{8}{\begin{aligned}
a_2 &\text{ durch } f, \\
a_3 &\text{ \textquotedbl} \ \varDelta, \\
a_4 &\text{ \textquotedbl} \ f^2, \\
a_5 &\text{ \textquotedbl} \ \varDelta f, \\
a_6 &\text{ \textquotedbl} \ \varDelta^2, f^3, \\
a_7 &\text{ \textquotedbl} \ C, \varDelta f^2
\end{aligned}}
